\section{An exact flux identity for unvisited near misses}
\label{sec:cf-flux}
The incidence load does not see a trajectory whose original guard is
witnessed only by clearance. We introduce a counting observable on the
\emph{same} physical flights, not a new collision map. Throughout this
section let
\[
 g_*=1-2R_+=3/50,\qquad \delta_*=g_*/4=3/200,
 \qquad 0<\delta\le\delta_*.
\]
The symbol $\delta_*$ here is a one-flight geometric constant, not the
outer microscopic width $\delta(B)$. All disks still have radius $R$.

For a regular outgoing collision $x$, write its flight as
$q(x)+tv(x)$, $0\le t\le\tau(x)$. Let $\mathcal Z_\delta(x)$ be the
set of lattice centers other than the two incident centers for which
\[
 0<\min_{0\le t\le\tau(x)}|q(x)+tv(x)-z|-R<\delta,
 \qquad N_\delta(x)=|\mathcal Z_\delta(x)|.
\]
This counts near misses, not collisions. In particular neither the
collision index, the physical roof, nor the section-visit count is
incremented by $N_\delta$. A center is counted once on a flight.
Tangencies may be assigned value zero on their collision-area null set;
no assertion about a positive pairing-endpoint trace follows from this
convention.

\begin{lemma}[An interior foot and uniform endpoint separation]
\label{lem:cf-foot}
Every center counted by $N_\delta$ has its perpendicular foot in the
open physical flight. If $t_z$ is the time of that foot, then
\begin{equation}\label{eq:cf-gap}
 t_z\ge g_*-\delta,\qquad \tau-t_z\ge g_*-\delta.
\end{equation}
The candidate centers, after translating the initial incident center
to zero, belong to the finite set
\[
 \{z\in\Lambda\setminus\{0\}:|z|\le\tau_++2R_++\delta_*\}.
\]
Write $N_*$ for its cardinality. This constant is independent of the
itinerary length and of all exact record labels.
\end{lemma}
\begin{proof}
At either incident point, the distance to a third center is at least
$1-R$, since distinct lattice centers have distance at least one.
Its distance to the third disk is therefore at least $1-2R\ge g_*$.
A minimum below $\delta\le g_*/4$ cannot occur at an endpoint. The
nearest point on a line segment is then its perpendicular foot, whose
distance from $z$ is at most $R+\delta$. By the reverse triangle
inequality its distance along the line from either endpoint is at
least $(1-R)-(R+\delta)\ge g_*-\delta$. Finally a foot is at distance
at most $R+\tau_+$ from the initial center. Adding its distance from
$z$ gives the finite-set bound. This proof uses the actual segment,
including both of its endpoints, rather than an infinite-line
substitute for the clearance event.
\end{proof}

Put $v_\vartheta=(\cos\vartheta,\sin\vartheta)$ and
$v_\vartheta^\perp=(-\sin\vartheta,\cos\vartheta)$. For a mark $z$
define the signed line coordinate
$b=(z-q)\cdot v_\vartheta^\perp$. Translate the marked center to zero.
For $\sigma\in\{-1,1\}$ and $0<u<\delta_*$ the oriented line with
$b=\sigma(R+u)$ has closest point
$-\sigma(R+u)v_\vartheta^\perp$. Follow this line backwards to its
last actual collision. Reduce that collision by a lattice translation
to the initial fundamental disk, and denote the resulting outgoing
state and marked-center displacement by
\[
 X_R(\vartheta,\sigma,u),\qquad Z_R(\vartheta,\sigma,u).
\]
These maps use straight free flight until an \emph{actual} disk is
met. There is no reflection at the closest point. The finite-horizon
assumption ensures actual collisions in both directions, apart from
null singular parameters on which these maps can be defined
arbitrarily. The interval from the preceding collision to the following
one remains the original physical flight.

\begin{theorem}[One-copy clearance flux and exact multiplicity]
\label{thm:cf-flux}
For every nonnegative Borel function $a$ of an outgoing state and a
marked-center displacement,
\begin{equation}\label{eq:cf-flux}
 \int\sum_{z\in\mathcal Z_\delta(x)}a(x,z)\,d\nu(x)
 =\frac{1}{4\pi R}\sum_{\sigma=\pm1}
   \int_0^{2\pi}\int_0^\delta
       a(X_R(\vartheta,\sigma,u),Z_R(\vartheta,\sigma,u))
                                  \,du\,d\vartheta.
\end{equation}
In particular,
\begin{equation}\label{eq:cf-exact-mean}
                         \int N_\delta\,d\nu=\delta/R.
\end{equation}
Both statements hold for every $R\in[R_-,R_+]$. The multiplicity on
the left is the number of physically distinct unvisited disks, not
an upper bound for the number of collision words.
\end{theorem}
\begin{proof}
We first verify the local measure, including its radius factor. Write
$q=R(\cos\alpha,\sin\alpha)$ and
$\vartheta=\alpha+\arcsin p$. At fixed $\vartheta$,
\[
 \left.\frac{\partial b}{\partial\alpha}\right|_\vartheta
     =-R\cos(\vartheta-\alpha),\qquad
 \frac{\partial\vartheta}{\partial p}=(1-p^2)^{-1/2}.
\]
Consequently
\begin{equation}\label{eq:cf-jacobian}
 \left|\det D_{(\alpha,p)}(\vartheta,b)\right|=R,
 \qquad d\nu=\frac{d\vartheta\,db}{4\pi R}.
\end{equation}
The two signs $b=\sigma(R+u)$ both have $|db|=du$.

There is also a global multiplicity assertion to prove. A marked
physical flight maps to its unique oriented line relative to its
marked center. Conversely, the closest point of every line just
described is in the free region: its distance to any other lattice
center is at least $1-R-u>R$. In fact the disk of radius $R+\delta_*$
around the marked center does not intersect any other actual disk,
since $2R+\delta_*<1$. The central disk is missed because the line
has impact magnitude greater than $R$. Following the line in both
directions therefore gives a physical flight whose incident disks
are different from the marked disk. Its foot is interior and its
clearance from that disk is $u$. This is inverse to the marked-flight
map, modulo the prescribed lattice translation. Multiple near misses
on one flight give different marked copies; they are neither lost nor
identified.

For measure-theoretic precision, partition the line parameters by the
actual backward incident disk and the forward incident disk, and
exclude their tangent boundaries. These partitions are Borel; on each
regular piece the disk-intersection formulas and their inverses are
smooth. Tangency boundaries satisfy the line equation $|b_z|=R$.
For any fixed alternative center their equation in $(\vartheta,u)$
is a graph in $u$ (with one of the two fixed signs), hence has planar
Lebesgue measure zero. There are countably many centers. The
corresponding singular collision sets are also null by
\eqref{eq:cf-jacobian} on regular charts. Exhaust these charts by
relatively compact subcharts and disjointize them. Applying the local
Jacobian formula to bounded functions on finite unions, then using
monotone convergence for both chart unions and $a\wedge n$, proves
\eqref{eq:cf-flux} without a prior integrability assumption. Setting
$a=1$ gives $2(2\pi)\delta/(4\pi R)=\delta/R$.
\end{proof}

\begin{corollary}[All physical depths with the original normalization]
\label{cor:cf-depth-mass}
For arbitrary $0<\delta_j\le\delta_*$, $0\le j<m$,
\[
 \int\sum_{j=0}^{m-1}N_{\delta_j}(T_R^jx)\,d\nu(x)
       =\frac1R\sum_{j=0}^{m-1}\delta_j.
\]
For every original exact event $\Pi$ the same integral with
$\one_\Pi\,d\nu_R^*$ is at most
$(Rc_*)^{-1}\sum_j\delta_j$. In particular, if $2\varepsilon\le\delta_*$
and $r=q_0^{1/4}$, then
\begin{equation}\label{eq:cf-all-depth-mass}
 \nu_R^*\!\left(\Pi\cap
   \{\exists j:\Delta_j<2\varepsilon r^{d'_j}\}\right)
       \le\frac{4\varepsilon}{R_-c_*(1-r)}.
\end{equation}
\end{corollary}
\begin{proof}
The collision measure $\nu$ is invariant under the original elastic
collision map. This also follows from the oriented-line Jacobian
\eqref{eq:cf-jacobian}, whose magnitude is the same at both incident
circles, and specular reflection, which preserves boundary flux.
Thus apply \eqref{eq:cf-exact-mean} at each $j$. The original initial
restriction has density $\one_{Y_R^*}/c_*\le1/c_*$ with respect to
$\nu$; further exact labels only decrease the positive integral.
Outside the null exact-tangency set, existence of a clearance witness
implies $N_{2\varepsilon r^{d'_j}}\ge1$. There are at most two flight
indices of each nonnegative depth, so
$\sum_j2\varepsilon r^{d'_j}\le4\varepsilon/(1-r)$. No independence
or return-conditioned invariant measure is introduced.
\end{proof}
The estimate \eqref{eq:cf-all-depth-mass} is an integrated mass
estimate. Multiplication of it by $m^2$ at fixed $\varepsilon(B)$
is not a proof about a target-roof density.
